\section{Ôn tập chương 3}
\subsection{Đề 1}
\caulc
\Opensolutionfile{ans}[ans/ans-De1-LC]
%%%==============EX_1============%%%
\begin{ex}%[2D3N1-3]
	Cho mẫu số liệu ghép nhóm có tứ phân vị thứ nhất, thứ hai, thứ ba lần lượt là $Q_1$, $Q_2$, $Q_3$. Khoảng tứ phân vị của mẫu số liệu ghép nhóm đó bằng
	\choice
	{$2Q_2$}
	{$Q_1-Q_3$}
	{\True $Q_3-Q_1$}
	{$Q_3+Q_1-Q_2$}
	\loigiai{
		Khoảng tứ phân vị của mẫu số liệu ghép nhóm đó là $\Delta_Q=Q_3-Q_1$.
	}
\end{ex}
%%%==============EX_2============%%%
\begin{ex}%[2D3N2-1]
	Một vườn thú ghi lại tuổi thọ (đơn vị: năm) của $20$ con hổ và thu được kết quả như sau
	\begin{center}
		\begin{tabular}{|c|c|c|c|c|c|} 
			\hline 
			Tuổi thọ & $\left[14;15\right)$ & $\left[15;16\right)$ & $\left[16;17\right)$ & $\left[17;18\right)$ & $\left[18;19\right)$ \\ 
			\hline 
			Số con hổ & $1$ & $3$ & $8$ & $6$ & $2$ \\ 
			\hline 
		\end{tabular}		
	\end{center}
	Khoảng biến thiên của mẫu số liệu ghép nhóm này là
	\choice
	{$3$}
	{$4$}
	{\True $5$}
	{$6$}
	\loigiai{
		Khoảng biến thiên $R=19-14=5$.
	}
\end{ex}
%%%==============EX_3============%%%
\begin{ex}%[2D3H1-3]
	Một vườn thú ghi lại tuổi thọ (đơn vị: năm) của $20$ con hổ và thu được kết quả như sau
	\begin{center}
		\begin{tabular}{|c|c|c|c|c|c|} 
			\hline 
			Tuổi thọ & $\left[14;15\right)$ & $\left[15;16\right)$ & $\left[16;17\right)$ & $\left[17;18\right)$ & $\left[18;19\right)$ \\ 
			\hline 
			Số con hổ & $1$ & $3$ & $8$ & $6$ & $2$ \\ 
			\hline 
		\end{tabular}		
	\end{center}
	Nhóm chứa tứ phân vị thứ nhất là
	\choice
	{$\left[14;15\right)$}
	{$\left[15;16\right)$}
	{\True $\left[16;17\right)$}
	{$\left[17;18\right)$}
	\loigiai{
		Cỡ mẫu là $1+3+8+6+2=20$.\\
		Gọi $x_1;x_2;\ldots;x_{20}$ là tuổi thọ của $20$ con hố được sắp xếp theo thứ tự tăng dần.\\
		Tứ phân vị thứ nhất của mẫu số liệu gốc là $\dfrac{x_5+x_6}{2}$.\\
		Mà $x_5;x_6$ đều thuộc nhóm $[16; 17)$ nên nhóm chứa tứ phân vị thứ nhất là $[16; 17)$.
	}
\end{ex}
%%%==============EX_4============%%%
\begin{ex}%[2D3H1-3]
	Một vườn thú ghi lại tuổi thọ (đơn vị: năm) của $20$ con hổ và thu được kết quả như sau
	\begin{center}
		\begin{tabular}{|c|c|c|c|c|c|} 
			\hline 
			Tuổi thọ & $\left[14;15\right)$ & $\left[15;16\right)$ & $\left[16;17\right)$ & $\left[17;18\right)$ & $\left[18;19\right)$ \\ 
			\hline 
			Số con hổ & $1$ & $3$ & $8$ & $6$ & $2$ \\ 
			\hline 
		\end{tabular}		
	\end{center}
	Nhóm chứa tứ phân vị thứ ba là
	\choice
	{$\left[15;16\right)$}
	{$\left[16;17\right)$}
	{\True $\left[17;18\right)$}
	{$\left[18;19\right)$}
	\loigiai{
		Tứ phân vị thứ ba của mẫu số liệu gốc là $\dfrac{x_{15}+x_{16}}{2}$.
		Mà $x_{15}$; $x_{16}$ đều thuộc nhóm $[17; 18)$. Do đó nhóm chứa tứ phân vị thứ ba là $[17; 18)$.
	}
\end{ex}
%%%==============EX_5============%%%
\begin{ex}%[2D3N1-1]
	Một vườn thú ghi lại tuổi thọ (đơn vị: năm) của $20$ con hổ và thu được kết quả như sau
	\begin{center}
		\begin{tabular}{|c|c|c|c|c|c|} 
			\hline 
			Tuổi thọ & $\left[14;15\right)$ & $\left[15;16\right)$ & $\left[16;17\right)$ & $\left[17;18\right)$ & $\left[18;19\right)$ \\ 
			\hline 
			Số con hổ & $1$ & $3$ & $8$ & $6$ & $2$ \\ 
			\hline 
		\end{tabular}		
	\end{center}
	Số đặc trưng nào không sử dụng thông tin của nhóm số liệu đầu tiên và nhóm số liệu cuối cùng.
	\choice
	{Khoảng biến thiên}
	{\True Khoảng tứ phân vị}
	{Phương sai}
	{Độ lệch chuẩn}
	\loigiai{
		Số đặc trưng không sử dụng thông tin của nhóm số liệu đầu tiên và nhóm số liệu cuối cùng là khoảng tứ phân vị.
	}
\end{ex}
%%%==============EX_6============%%%
\begin{ex}%[2D3N1-2]
	Một vườn thú ghi lại tuổi thọ (đơn vị: năm) của $20$ con hổ và thu được kết quả như sau
	\begin{center}
		\begin{tabular}{|c|c|c|c|c|c|} 
			\hline 
			Tuổi thọ & $\left[14;15\right)$ & $\left[15;16\right)$ & $\left[16;17\right)$ & $\left[17;18\right)$ & $\left[18;19\right)$ \\ 
			\hline 
			Số con hổ & $1$ & $3$ & $8$ & $6$ & $2$ \\ 
			\hline 
		\end{tabular}		
	\end{center}
	Nếu thay tất cả các tần số trong mẫu số liệu ghép nhóm trên bằng $4$ thì số đặc trưng nào sau đây không thay đổi?
	\choice
	{\True Khoảng biến thiên}
	{Khoảng tứ phân vị}
	{Phương sai}
	{Độ lệch chuẩn}
	\loigiai{
		Khoảng biến thiên sẽ không thay đổi nếu thay tất cả các tần số trong mẫu số liệu ghép nhóm trên bằng $4$.
	}
\end{ex}
%%%=============EX_7=============%%%
\begin{ex}%[2D3H2-2]
	Cho mẫu số liệu thống kê $\{1; 2; 3; 4; 5; 6; 7; 8; 9\}$. Tính (gần đúng) độ lệch chuẩn của mẫu số liệu trên.
	\choice
	{$2{,}45$}
	{\True $2{,}58$}
	{$6{,}67$}
	{$6{,}0$}
	\loigiai{
		Ta có
		\begin{itemize}
			\item Số trung bình $\overline{x}=\dfrac{1}{9}\left(1+2+3+4+5+6+7+8+9\right)=5$.
			\item Phương sai $s^2=\dfrac{1}{9}\left[(5-1)^2+(5-2)^2+\cdots+(5-9)^2\right]=\dfrac{20}{3}$.
			\item Độ lệch chuẩn $s=\sqrt{s^2}\approx 2{,}58$.
		\end{itemize}
	}
\end{ex}

%%%=============EX_8=============%%%
\begin{ex}%[2D3H2-2]
	Cho mẫu số liệu thống kê $\{2,4,6,8,10\}$. Phương sai của mẫu số liệu trên bằng bao nhiêu?
	\choice
	{$6$}
	{\True $8$}
	{$10$}
	{$40$}
	\loigiai{
		Ta có
		\begin{itemize}
			\item Số trung bình $\overline{x}=\dfrac{1}{5}\left(2+4+6+8+10\right)=6$.
			\item Phương sai $s^2=\dfrac{1}{5}\left[(6-2)^2+(6-4)^2+(6-6)^2+(6-8)^2+(6-10)^2\right]=8$.
		\end{itemize}
	}
\end{ex}

%%%=============EX_9=============%%%
\begin{ex}%[2D3H2-2]
	Cho dãy số liệu thống kê $1,2,3,4,5,6,7$. Phương sai của mẫu số liệu thống kê đã cho là
	\choice
	{$1$}
	{$2$}
	{$3$}
	{\True $4$}
	\loigiai{
		Ta có
		\begin{itemize}
			\item Số trung bình $\overline{x}=\dfrac{1}{7}\left(1+2+3+4+5+6+7\right)=4$.
			\item Phương sai $s^2=\dfrac{1}{7}\left[(4-1)^2+(4-2)^2+\cdots+(4-7)^2\right]=4$.
		\end{itemize}
	}
\end{ex}

%%%=============EX_10=============%%%
\begin{ex}%[2D3H2-2]
	Cho mẫu số liệu $10,8,6,2,4$. Độ lệch chuẩn của mẫu là
	\choice
	{$8$}
	{$2{,}4$}
	{$6$}
	{\True $2{,}8$}
	\loigiai{
		Ta có
		\begin{itemize}
			\item Số trung bình $\overline{x}=\dfrac{1}{5}\left(2+4+6+8+10\right)=6$.
			\item Phương sai $s^2=\dfrac{1}{5}\left[(6-2)^2+(6-4)^2+\cdots+(6-10)^2\right]=8$.
			\item Độ lệch chuẩn $s=\sqrt{s^2}\approx 2{,}83$.
		\end{itemize}
	}
\end{ex}

%%%=============EX_11=============%%%
\begin{ex}%[2D3H2-2]
	Bảng số liệu sau cho biêt thời gian làm bài tính bằng phút của 50 học sinh.
	\begin{center}
		\begin{tabular}{|c|c|c|c|c|c|c|c|c|c|c|c|}
			\hline
			Thời gian & 3 & 4 & 5 & 6 & 7 & 8 & 9 & 10 & 11 & 12 & \\
			\hline
			Tần số (n) & 1 & 3 & 4 & 7 & 8 & 9 & 8 & 5 & 3 & 2 & $\mathrm{~N}=50$ \\
			\hline
		\end{tabular}
	\end{center}
	
	Tìm độ lệch chuẩn của mẫu số liệu thống kê trên.
	\choice
	{$\delta \approx 2{,}15$}
	{$\delta \approx 2{,}14$}
	{$\delta \approx 2{,}16$}
	{\True $\delta \approx 2{,}13$}
	\loigiai{
		Ta có
		\begin{itemize}
			\item Thời gian làm bài trung bình \[\overline{x}=\dfrac{1}{50}\left(3\cdot 1+4\cdot 3+5\cdot 4+6\cdot 7+7\cdot 7+8\cdot 9+9\cdot 8+10\cdot 5+11\cdot 3+12\cdot 2\right)=7{,}68.\]
			\item Phương sai $\delta^2=\dfrac{1}{5}\left[1\cdot (7{,}68-3)^2+3\cdot (7{,}68-4)^2+\cdots+2\cdot (7{,}68-12)^2\right]\approx 4{,}54$.
			\item Độ lệch chuẩn $\delta=\sqrt{\delta^2}\approx 2{,}13$.
		\end{itemize}
	}
\end{ex}

%%%=============EX_12=============%%%
\begin{ex}%[2D3H2-2]
	Sản lượng lúa (đơn vị ha) của $40$ thửa ruộng có cùng diện tích được trình bày trong bảng số liệu sau
	\begin{center}
		\begin{tabular}{|c|c|c|c|c|c|c|}
			\hline
			Sản lượng & 20 & 21 & 22 & 23 & 24 & \\
			\hline
			Tần số & 5 & 8 & 11 & 10 & 6 & $N=40$ \\
			\hline
		\end{tabular}
	\end{center}
	
	Tính phương sai của bảng số liệu.
	\choice
	{$1{,}55$}
	{$1{,}53$}
	{$1{,}52$}
	{\True $1{,}54$}
	\loigiai{
		Ta có
		\begin{itemize}
			\item Sản lượng lúa trung bình \[\overline{x}=\dfrac{1}{40}\left(20\cdot 5+21\cdot 8+22\cdot 11+23\cdot 10+24\cdot 6\right)=22{,}1.\]
			\item Phương sai $\delta^2=\dfrac{1}{40}\left[5\cdot (22{,}1-20)^2+8\cdot (22{,}1-21)^2+\cdots+6\cdot (22{,}1-24)^2\right]\approx 1{,}54$.
		\end{itemize}
	}
\end{ex}
\Closesolutionfile{ans}
\indapan{6}{ans/ans-De1-LC}

\cauds
\Opensolutionfile{ans}[ans/ans-De1-DS]
\begin{ex}%[2D3V1-4]
Một trường trung học phổ thông thống kê lại chiều cao của một số học sinh và được tổng hợp ở bảng bên dưới.
\begin{longtable}{|c|c|c|c|}
\hline
\textbf{Nhóm}&\textbf{Chiều cao} ($\mathrm{cm}$)&\textbf{Giá trị đại diện} ($\mathrm{cm}$)&\textbf{Số học sinh}\\
\hline
$1$&$[150;153)$&$151{,}5$&$7$\\
\hline
$2$&$[153;156)$&$154{,}5$&$13$\\
\hline
$3$&$[156;159)$&$157{,}5$&$40$\\
\hline
$4$&$[159;162)$&$160{,}5$&$21$\\
\hline
$5$&$[162;165)$&$163{,}5$&$13$\\
\hline
$6$&$[165;168)$&$166{,}5$&$6$\\
\hline
\end{longtable}
\choiceTFt
{\True Khoảng biến thiên về chiều cao của các học sinh bằng $18\mathrm{~cm}$}
{\True Giá trị tứ phân vị thứ nhất của mẫu số liệu trên là $Q_1=156{,}375$}
{Giá trị tứ phân vị thứ ba của mẫu số liệu trên là $Q_3=160{,}5$}
{\True Mẫu số liệu ghép nhóm trên không có giá trị ngoại lệ}
\loigiai{
\begin{itemchoice}
	\itemch $R=x_n-x_1=168-150=18$.
	\itemch Ta có cỡ mẫu $n=100$ nên tứ phân vị thứ nhất của mẫu số liệu là\\ $\dfrac{1}{2}(x_{25}+x_{26})\in [156;159)$.\\
	Do đó $Q_1=156+\dfrac{\dfrac{ 100}{4}-(7+13)}{40}(159-156)=\dfrac{1\,251}{8}=156{,}375$.
	\itemch Tứ phân vị thứ ba của mẫu số liệu là $\dfrac{1}{2}(x_{75}+x_{76})\in [159;162)$.\\
	Do đó  $Q_k=u_m+\dfrac{\dfrac{kn}{4}-C}{n_m}(u_{m+1}-u_{m})$\\$\Rightarrow Q_3=159+\dfrac{\dfrac{3\cdot 100}{4}-(7+13+40)}{21}(162-159)=\dfrac{1\,128}{7}$.\\
		\itemch 
Khoảng tứ phân vị $\Delta Q=Q_3-Q_1=\dfrac{267}{56}$.\\
Ta có $\heva{&Q_1-1{,}5\Delta Q=\dfrac{1\,251}{8}-1{,}5\cdot \dfrac{267}{56}=\dfrac{16\,713}{112}\approx 149{,}2\\&Q_3+1{,}5\Delta Q=\dfrac{1\,128}{7}+1{,}5\cdot \dfrac{267}{56}=\dfrac{18\,849}{112}\approx 168{,}3.}$\\
Vậy mẫu số liệu không có giá trị ngoại lệ.
\end{itemchoice}
}
\end{ex}

\begin{ex}%[2D3V1-4]
	Điểm thi đánh giá năng lực của một trường đại học vào năm $2024$ có biểu đồ như hình bên dưới.
	\begin{center}
		\begin{tikzpicture}[scale=1, font=\footnotesize, line join=round, line cap=round, >=stealth]
			\coordinate [label= left: $0$](O) at (0,0) ;
			\coordinate [label= left: Số học sinh](y) at (0,8) ;
			\coordinate [label= below left: Điểm](x) at (14.3,0) ;
			\foreach \i in {1,2,3,4,5,6,7}
			{\draw[-,color=gray!40] (0,\i)--(14,\i); }
			\draw (0,1)node[left]{$20$}(0,2)node[left]{$40$}(0,3)node[left]{$60$}(0,4)node[left]{$80$}(0,5)node[left]{$100$}(0,6)node[left]{$120$}(0,7)node[left]{$140$};
			%--------------------------------------------
			\foreach \i/\j/\a in 
			{0.25/0/0,0.75/0/0,1.25/0/0,1.75/0/0,2.25/0/0,2.75/0/0,
				3.25/0/0,3.75/0.05/1,4.25/0.4/8,4.75/1.2/24,5.25/2.7/54,
				5.75/4.75/95,6.25/4.75/95,6.75/6.65/133,7.25/6.1/122,	7.75/5.2/104,8.25/3.1/62,8.75/2.75/55,9.25/1.05/21,9.75/0.6/12,
				10.25/0/0,10.75/0/0,11.25/0/0,11.75/0/0}{
				\draw[] ($(\i,\j)+(0,0.2)$) node[scale=0.9]{$\a$};
			}
			%-----------------------------------------------
			\foreach \i/\j in 
			{0/0,0.5/0,1/0,1.5/0,2/0,2.5/0,3/0,
				3.5/0.05,4/0.4,4.5/1.2,5/2.7,5.5/4.75,6/4.75,
				6.5/6.65,7/6.1,7.5/5.2,8/3.1,8.5/2.75,9/1.05,
				9.5/0.6,10/0,10.5/0,11/0,11.5/0}{
				\fill[blue!20] (\i,0) rectangle+ (0.5,\j);
				\draw[red,thick] (\i,0) rectangle+ (0.5,\j);
			}
			\foreach \i/\j in 
			{0.25/[0-50),0.75/[50-100),1.25/[100-150),1.75/[150-200),2.25/[200-250),2.75/[250-300),3.25/[300-350),3.75/[350-400),4.25/[400-450),4.75/[450-500),5.25/[500-550),5.75/[550-600),6.25/[600-650),6.75/[650-700),7.25/[700-750),7.75/[750-800),8.25/[800-850),8.75/[850-900),9.25/[900-950),9.75/[950-1000),10.25/[1000-1050),10.75/[1050-1100),11.25/[1100-1150),11.75/[1150-1200)}{
				\draw($(\i,0)+(0,-1)$) node[rotate=90]{\j};
			}	
			\foreach \i/\j in {}{\draw (\i,\j) node[above]{\j};
			}			
			\draw[->,very thick] (O)--(x);
			\draw[->,very thick] (O)--(y);
		\end{tikzpicture}
	\end{center}	
\choiceTFt
{Có $140$ học sinh tham gia kì thi đánh giá năng lực}
{\True Giá trị đại diện cho nhóm chứa tứ phân vị thứ ba của mẫu số liệu ghép nhóm là $775$}
{Giá trị tứ phân vị thứ nhất của mẫu số liệu ghép nhóm là $Q_1=625$}
{\True Mẫu số liệu ghép nhóm trên có giá trị ngoại lệ $x\in [350;400)$}
\loigiai{
\begin{itemchoice}
\itemch Số học sinh tham gia\\ $n=1+8+24+54+95+95+133+122+104+62+55+21+12=786$.
\itemch Tứ phân vị thứ ba của mẫu số liệu là $x_{589}\in [750;800)$.\\
Giá trị đại diện cho nhóm chứa tứ phân vị thứ ba là $\dfrac{750+800}{2}=775$.
\itemch Ta có $n=786$ nên tứ phân vị thứ nhất của mẫu số liệu là $x_{196}\in [600;650)$.\\
Do đó giá trị của tứ phân vị thứ nhất là\\ $Q_1=600+\dfrac{\dfrac{786}{4}-(1+8+24+54+95)}{95}\cdot (650-600)=\dfrac{11\,545}{19}$.
\itemch Ta có $n=786$ nên tứ phân vị thứ ba của mẫu số liệu là $x_{589}\in [750;800)$.\\
Do đó giá trị của tứ phân vị thứ ba là\\ $Q_3=750+\dfrac{\dfrac{3\cdot786}{4}-(1+8+24+54+95+133+122)}{104}\cdot (800-750)=\dfrac{80\,875}{104}$.\\
Khoảng tứ phân vị $\Delta Q=Q_3-Q_1=\dfrac{80\,875}{104}-\dfrac{11\,545}{19}\approx 170{,}01$.\\
Ta có: $\heva{&Q_3+1{,}5\Delta Q\approx 1\,032{,}66\\&Q_1-1{,}5\Delta Q\approx 352{,}62.}$\\
Vậy mẫu số liệu có giá trị ngoại lệ $x<352{,}62$.
\end{itemchoice}
}
\end{ex}


\begin{ex}%[2D3V2-3]
Hình bên dưới là biểu đồ biểu diễn nhiệt độ trung bình hằng tháng của hai địa phương $Y$ và $Z$.
\begin{center}
	\begin{tikzpicture}[scale=1, line join=round, line cap=round, >=stealth]
		\coordinate [label= left: $0$](O) at (0,0) ;
		\coordinate [label= left: $^\circ\mathrm{C}$](y) at (0,9) ;
		\coordinate [label= below: Tháng](x) at (14,0) ;
\draw[line width=1pt,color=gray!50] (0,0) grid (14,9);
\foreach \i in {0.2,0.4,0.6,...,9}
{\draw[-,color=gray!10] (0,\i)--(14,\i); }
%--------------------------------------------
\draw (0,1)node[left]{$5$}(0,2)node[left]{$10$}(0,3)node[left]{$15$}(0,4)node[left]{$20$}(0,5)node[left]{$25$}(0,6)node[left]{$30$}(0,7)node[left]{$35$}(0,8)node[left]{$40$};
\draw (1,0)node[below]{$1$}(2,0)node[below]{$2$}(3,0)node[below]{$3$}(4,0)node[below]{$4$}(5,0)node[below]{$5$}(6,0)node[below]{$6$}(7,0)node[below]{$7$}(8,0)node[below]{$8$}(9,0)node[below]{$9$}(10,0)node[below]{$10$}(11,0)node[below]{$11$}(12,0)node[below]{$12$};
%--------------------------------------------
	\foreach \i/\j\a in 
{1/2.2/0.2,2/2.4/0.8,3/3.2/0.6,4/3.8/0.4,5/4.2/1.2,6/5.4/0.6,7/6/1.2,8/7.2/0.5,9/7.7/-1.2,10/6.5/-0.5,11/6/-1.7}{
	\fill[blue] (\i,\j) circle (3pt)(12,4.3)circle (3pt);
	\draw[blue,line width=1.5pt] (\i,\j)--($(\i+1,\j+\a)$);
}		
%--------------------------------------------		
\foreach \i/\j\a in 
{1/4.6/0.9,2/5.5/0.5,3/6/0.2,4/6.2/0.4,5/6.6/0.4,6/7/0.3,7/7.3/-1.1,8/6.2/-0.4,9/5.8/0,10/5.8/-0.4,11/5.4/-1.6}{
	\fill[red] (\i,\j) circle (3pt)(12,3.8)circle (3pt);
	\draw[red,line width=1.5pt] (\i,\j)--($(\i+1,\j+\a)$);
}				
%--------------------------------------------
\foreach \i/\j in {}{\draw (\i,\j) node[above]{\j};}			
		\draw[->,very thick] (O)--(x);
		\draw[->,very thick] (O)--(y);
\path (9.7,7) node[above]{$(Y)$};
\path (9.5,5.1) node[above]{$(Z)$};	
\path (7,10)node[below]{\textbf{Biểu đồ nhiệt trung bình hằng tháng của hai địa phương $Y$ và $Z$}};
	\end{tikzpicture}
\end{center}
\choiceTFt
{Với độ dài các nhóm là $5$ và đầu mút phải của nhóm cuối cùng là $40$ thì bảng số số liệu ghép nhóm về nhiệt độ trung bình của hai địa phương $Y,Z$ đều có $6$ nhóm}
{\True Giá trị tứ phân vị thứ nhất của địa phương $Y$ bằng $17{,}5$}
{\True Giá trị tứ phân vị thứ ba của địa phương $Z$ bằng $32{,}5$}
{Nếu so sánh theo độ lệch chuẩn thì nhiệt độ trung bình trong một năm của địa phương $Y$ ít phân tán hơn địa phương $Z$}
\loigiai{
Ta có bảng tần số ghép nhóm về nhiệt độ trung bình của hai địa phương $Y$ và $Z$ như sau
\begin{longtable}{|c|c|c|c|c|c|c|}
\hline
\textbf{Nhiệt độ trung bình} $^\circ C$&$[10;15)$&$[15;20)$&$[20;25)$&$[25;30)$&$[30;35)$&$[35;40)$\\
\hline
\textbf{Địa phương} $Y$&$2$&$2$&$2$&$1$&$3$&$2$\\
\hline
\textbf{Địa phương} $Z$&$0$&$1$&$1$&$5$&$4$&$1$\\
\hline
\end{longtable}\vspace{-5mm}
\begin{itemchoice}
\itemch Dựa vào bảng tần số ghép nhóm, ta có địa phương $Y$ chia thành $6$ nhóm và địa phương $Z$ chia thành $5$ nhóm.
\itemch Cỡ mẫu $n=12$ nên tứ phân vị thứ nhất của địa phương $Y$ là $\dfrac{1}{2}\left(x_3+x_4\right)\in [15;20)$.\\
Giá trị tứ phân vị thứ nhất của địa phương $Y$ $$Q_1=15+\dfrac{\dfrac{12}{4}-2}{2}\cdot (20-15)=17{,}5.$$
\itemch Cỡ mẫu $n=12$ nên tứ phân vị thứ ba của địa phương $Z$ là $\dfrac{1}{2}\left(x_9+x_{10}\right)\in [30;35)$.\\
Giá trị tứ phân vị thứ ba của địa phương $Z$ là $$Q_3=30+\dfrac{\dfrac{3\cdot 12}{4}-(1+1+5)}{4}\cdot (35-30)=32{,}5.$$
\itemch Nhiệt độ trung bình của địa phương $Y$
$$\overline{x}_Y=\dfrac{1}{12}\left(12{,}5\cdot 2+17{,}5\cdot 2+22{,}5\cdot 2+27{,}5\cdot 1+32{,}5\cdot 3+37{,}5\cdot2\right)\approx 25{,}4.$$
Phương sai về nhiệt độ trung bình của địa phương $Y$
$$S^2_Y=\dfrac{1}{2}\left(12{,}5^2\cdot 2+17{,}5^2\cdot 2+22{,}5^2\cdot 2+27{,}5^2\cdot 1+32{,}5^2\cdot 3+37{,}5^2\cdot2 \right)-25{,}42^2\approx76{,}91.$$
Độ lệch chuẩn về nhiệt độ trung bình của địa phương $Y$ là $S_Y=\sqrt{S^2_Y}\approx 8{,}77$.\\
Nhiệt độ trung bình của địa phương $Z$
$$\overline{x}_Z=\dfrac{1}{12}\left(17{,}5\cdot 1+22{,}5\cdot 1+27{,}5\cdot 5+32{,}5\cdot 4+37{,}5\cdot1\right)= 28{,}75. $$
Phương sai về nhiệt độ trung bình của địa phương $Z$
$$S^2_Z=\dfrac{1}{2}\left(17{,}5\cdot 1+22{,}5^2\cdot 1+27{,}5^2\cdot 5+32{,}5^2\cdot 4+37{,}5^2\cdot 1 \right)-28{,}75^2\approx25{,}52.$$
Độ lệch chuẩn về nhiệt độ trung bình của địa phương $Z$ là $S_Z=\sqrt{S^2_X}\approx 5{,}05$.\\
Từ kết quả trên ta có $S_Y>S_Z$.
\end{itemchoice}
}
\end{ex}


\begin{ex}%[2D3V2-3]
Thống kê lại thu nhập trong một tháng của nhân viên hai công ty $A$ và $B$ (đơn vị: triệu đồng) được thể hiện trong biểu đồ dưới đây.
\begin{center}
\begin{tikzpicture}[scale=1.2, font=\footnotesize, line join=round, line cap=round, >=stealth]
\coordinate [label= left: $0$](O) at (0,0) ;
\coordinate [label= left: (Số nhân viên)](y) at (0,6) ;
\coordinate [label= below: (triệu đồng)](x) at (10,0) ;
\foreach \i in {1,2,3,4,5}
{\draw[-,color=gray!40] (0,\i)--(10,\i); }
\draw (0,1)node[left]{$10$}(0,2)node[left]{$20$}(0,3)node[left]{$30$}(0,4)node[left]{$40$}(0,5)node[left]{$50$};
%--------------------------------------------
\foreach \i/\j/\a in 
{0.4/1.2/12,1.2/1.6/16,2/4.5/45,2.8/4/40,3.6/1.6/16,4.4/2.9/29,
			5.2/1.4/14,6/1.5/15,6.8/0.7/7,7.6/1.2/12}{
\draw[] ($(\i,\j)+(0,0.2)$) node[scale=0.8]{$\a$};
}
%-----------------------------------------------
\foreach \i/\j in 
{0/1.2,1.6/4.5,3.2/1.6,4.8/1.4,6.4/0.7}{
\fill[blue!50] (\i,0) rectangle+ (0.8,\j);
\draw[thick] (\i,0) rectangle+ (0.8,\j);
}
	%-----------------------------------------------	
\foreach \i/\j in 
{0.8/1.6,2.4/4,4/2.9,5.6/1.5,7.2/1.2}{
\fill[pattern=north east lines] (\i,0) rectangle+ (0.8,\j);
\draw[thick] (\i,0) rectangle+ (0.8,\j);
}	
%-----------------------------------------------
\foreach \i/\j in 
{0.4/[5;9),2/[9;13),3.6/[13;17),5.2/[17;21),6.8/[21;25)}{
\draw($(\i,0)+(0.5,0)$) node[below]{\j};
}	
\foreach \i/\j in {}{\draw (\i,\j) node[above]{\j};
}			
\draw[->,very thick] (O)--(x);
\draw[->,very thick] (O)--(y);
\fill[blue!50,draw] (7.5,4.2) rectangle (8,4.8);
\fill[pattern=north east lines,draw] (7.5,3.2) rectangle (8,3.8);
\draw (8,4.5) node [right]{Công ty $A$};
\draw (8,3.5) node [right]{Công ty $B$};
\draw (5,5.5) node [above]{\textbf{Thu nhập trong một tháng của nhân viên hai công ty $A$ và $B$}};
	\end{tikzpicture}
\end{center}
\choiceTFt
{\True Có $14$ nhân viên của công ty $A$ thu nhập từ $17$ triệu đồng đến $21$ triệu đồng trong một tháng}
{Thu nhập trung bình mỗi tháng của nhân viên công ty $A$ cao hơn nhân viên công ty $B$}
{\True Nếu so sánh về phương sai thì thu nhập mỗi tháng của nhân viên công ty $A$ ít phân tán hơn nhân viên công ty $B$}
{\True Nếu so sánh về khoảng tứ phân vị thì thu nhập trung bình mỗi tháng của công ty $B$ đồng đều hơn công ty $A$}
\loigiai{
Ta có bảng tần số ghép nhóm về thu nhập trong một tháng của nhân viên hai công ty $A$ và $B$ như sau
\begin{longtable}{|c|c|c|c|c|c|}
\hline
\textbf{Thu nhập hàng tháng (triệu đồng)}&$[5;9)$&$[9;13)$&$[13;17)$&$[17;21)$&$[21;25)$\\
\hline
\textbf{Công ty $A$ (số nhân viên)}&$12$&$45$&$16$&$14$&$7$\\
\hline
\textbf{Công ty $B$ (số nhân viên)}&$16$&$40$&$29$&$15$&$12$\\
\hline
\end{longtable}\vspace{-5mm}
\begin{itemchoice}
\itemch Dựa vào bảng tần số ghép nhóm ở trên ta thấy có $14$ nhân viên của công ty $A$ thu nhập từ $17$ triệu đồng đến $21$ triệu đồng trong một tháng.
\itemch Số nhân viên công ty $A$ là $n=94$.\\
Thu nhập trung bình mỗi tháng của nhân viên công ty $A$
$$\overline{x}_A=\dfrac{1}{94}\left(7\cdot 12+11\cdot 45+15\cdot 16+19\cdot 14+23\cdot 7 \right)\approx13{,}26~\text{(triệu đồng).}$$
Số nhân viên công ty $B$ là $n=112$.\\
Thu nhập trung bình mỗi tháng của nhân viên công ty $B$
$$\overline{x}_B=\dfrac{1}{112}\left(7\cdot 16+11\cdot 40+15\cdot 29+19\cdot 15+23\cdot 12 \right)\approx 13{,}82~\text{(triệu đồng).}$$
Từ kết quả trên, ta có $\overline{x}_A<\overline{x}_B$.
\itemch 
Phương sai về thu nhập trung bình một tháng của nhân viên công ty $A$ $$S^2_A=\dfrac{1}{94}\left(7^2\cdot 12+11^2\cdot 45+15^2\cdot 16+19^2\cdot 14+23^\cdot 7 \right)-13{,}26^2\approx 19{,}93.$$
Phương sai về thu nhập trung bình một tháng của nhân viên công ty $B$ $$S^2_B=\dfrac{1}{112}\left(7^2\cdot 16+11^2\cdot 40+15^2\cdot 29+19^2\cdot 15+23^2\cdot 12 \right)-13{,}82^2\approx 22{,}47.$$
Từ kết quả trên, ta có $S^2_A<S^2_B$.
\itemch Tứ phân vị thứ nhất của nhân viên công ty $A$ là $x_{23}\in [9;13)$	nên $$Q_1=9+\dfrac{\dfrac{94}{4}-12}{45}(13-9)=\dfrac{451}{45}.$$
Tứ phân vị thứ ba của nhân viên công ty $A$ là $x_{71}\in [13;17)$	nên $$Q_3=13+\dfrac{\dfrac{4\cdot94}{4}-(12+45)}{16}(17-13)=\dfrac{89}{4}.$$
Khoảng tứ phân vị về thu nhập một tháng của nhân viên công ty $A$ $$\Delta Q_A=Q_3-Q_1=\dfrac{2\,201}{180}\approx 12{,}23.$$
Tứ phân vị thứ nhất của nhân viên công ty $B$ là $\dfrac{1}{2}\left( x_{28}+x_{29}\right) \in [9;13)$	nên $$Q_1=9+\dfrac{\dfrac{112}{4}-16}{40}(13-9)=\dfrac{51}{5}.$$
Tứ phân vị thứ ba của nhân viên công ty $B$ là $\dfrac{1}{2}\left( x_{84}+x_{85}\right) \in [13;17)$	nên $$Q_3=13+\dfrac{\dfrac{4\cdot112}{4}-(12+45)}{16}(17-13)=\dfrac{601}{29}.$$
Khoảng tứ phân vị về thu nhập một tháng của nhân viên công ty $B$ $$\Delta Q_B=Q_3-Q_1=\dfrac{1\,526}{145}\approx 10{,}52.$$	
Từ kết quả trên, ta có $\Delta Q_A>\Delta Q_B$.
\end{itemchoice}
}
\end{ex}

\Closesolutionfile{ans}
\indapan{3}{ans/ans-De1-DS}




\caukq
\Opensolutionfile{ans}[ans/ans-De1-KQ]
%%%============EX_1==============%%%
\begin{ex}%[2D3V1-2]
	Bạn Trang thống kê lại chiều cao (đơn vị: $\mathrm{cm}$) của các bạn học sinh nữ lớp 12C và lớp 12D ở bảng sau.
	\begin{center}
		\begin{tabular}{|c|c|c|c|c|c|c|}
			\hline
			Lớp Chiều cao & {$[155; 160)$} & {$[160; 165)$} & {$[165; 170)$} & {$[170; 175)$} & {$[175; 180)$} & {$[180; 185)$} \\
			\hline
			12C & 2 & 7 & 12 & 3 & 0 & 1 \\
			\hline
			12D & 5 & 9 & 8 & 2 & 1 & 0 \\
			\hline
		\end{tabular}
	\end{center}
	Gọi $\Delta_{Q}$ và $\Delta'_{Q}$ lần lượt là độ phân tán của nửa giữa hai mẫu số liệu chiều cao của các học sinh nữ lớp $12$C và $12$D. Tính giá trị của biểu thức $\Delta'_{Q}-\Delta_{Q}$ (kết quả làm tròn đến hàng phần trăm). 
	\shortans[]{$1{,}24$}
	\loigiai{
		\begin{itemize}
			\item Xét ở lớp $12$C, ta có cỡ mẫu $n=25$.\\
			Gọi $x_1; x_2; \ldots; x_{25}$ là mẫu số liệu gốc gồm chiều cao của $25$ bạn nữ lớp $12$C. Ta có 
			\begin{itemize}
				\item $x_1, x_2 \in[155; 160)$; 
				\item $x_3, \ldots, x_9 \in[160; 165)$; 
				\item $x_{10}, \ldots, x_{21} \in[165; 170)$; 
				\item $x_{22}, x_{23}, x_{24} \in[170; 175)$; 
				\item $x_{25} \in$ $[180; 185)$.
				\item Tứ phân vị thứ nhất của mẫu số liệu lớp $12$C là $\dfrac{x_6+x_7}{2} \in[160; 165)$. \\
				Do đó, tứ phân vị thứ nhất của mẫu số liệu lớp $12$C là
				\[Q_1=160+\dfrac{\dfrac{25}{4}-2}{7} \cdot(165-160)=\dfrac{4\,565}{28}.\]
				\item Tứ phân vị thứ ba của mẫu số liệu lớp $12$C là $\dfrac{x_{19}+x_{20}}{2} \in[165; 170)$.\\
				Do đó, tứ phân vị thứ ba của mẫu số liệu lớp $12C$ là \[Q_3=165+\dfrac{\dfrac{3 \cdot 25}{4}-(2+7)}{12} \cdot(170-165)=\dfrac{2\,705}{16}.\]
				Vậy khoảng tứ phân vị của mẫu số liệu lớp $12C$ là \[\Delta_Q=Q_3-Q_1=\dfrac{675}{112} \approx 6{,}03.\]
			\end{itemize}
			\item Xét ở lớp $12$D, ta có cỡ mẫu $n=25$.\\
			Gọi $x_1; x_2; \ldots; x_{25}$ là mẫu số liệu gốc gồm chiều cao của $25$ bạn nữ lớp $12$D. Ta có
			\begin{itemize}
				\item $x_1, \ldots, x_5 \in[155; 160)$; 
				\item $x_6, \ldots, x_{14} \in[160; 165)$; 
				\item $x_{15}, \ldots, x_{22} \in[165; 170)$; 
				\item $x_{23}, x_{24} \in[170; 175)$; 
				\item $x_{25} \in$ $[175; 180)$.
				\item Tứ phân vị thứ nhất của mẫu số liệu lớp $12$C là $\dfrac{x_6+x_7}{2} \in[160; 165)$.\\
				Do đó, tứ phân vị thứ nhất của mẫu số liệu lớp $12$C là
				\[Q_1=160+\dfrac{\dfrac{25}{4}-5}{9} \cdot(165-160)=\dfrac{5\,785}{36}.\]
				\item Tứ phân vị thứ ba của mẫu số liệu lớp $12$C là $\dfrac{x_{19}+x_{20}}{2} \in[165; 170)$.\\
				Do đó, tứ phân vị thứ ba của mẫu số liệu lớp $12C$ là
				\[Q_3=165+\dfrac{\dfrac{3 \cdot 25}{4}-(5+9)}{8} \cdot(170-165)=\dfrac{5\,375}{32}.\]				
			\end{itemize}
			Vậy khoảng tứ phân vị của mẫu số liệu lớp $12C$ là $\Delta'_Q=Q_3-Q_1 \approx 7{,}27$.\\
			Vậy giá trị của biểu thức $\Delta'_{Q}-\Delta_{Q}=7{,}27-6{,}03=1{,}24$.
		\end{itemize}
	}
\end{ex}
%%%==============EX_2============%%%
\begin{ex}%[2D3V1-2]
	Hằng ngày ông Thắng đều đi xe buýt từ nhà đến cơ quan. Dưới đây là bảng thống kê thời gian của $100$ lần ông Thắng đi xe buýt từ nhà đến cơ quan.
	\begin{center}
		\begin{tabular}{|c|c|c|c|c|c|c|}
			\hline
			Thời gian (phút) & {$[15; 18)$} & {$[18; 21)$} & {$[21; 24)$} & {$[24; 27)$} & {$[27; 30)$} & {$[30; 33)$} \\
			\hline
			Số lượt & 22 & 38 & 27 & 8 & 4 & 1 \\
			\hline
		\end{tabular}
	\end{center}
	Hāy tìm khoảng tứ phân vị của mẫu số liệu ghép nhóm ở dữ liệu trên sau khi đã loại bỏ các giá trị ngoại lệ (kết quả làm tròn đến hàng phần trăm). 
	\shortans[]{$4{,}37$}
	\loigiai{
		Ta có $Q_3+1,5\Delta_Q=\dfrac{6\,683}{228} < 30$ và $Q_1-1,5\Delta_Q=\dfrac{881}{76} \approx 11{,}59$.\\
		Trong lần duy nhất ông Thắng đi hết hơn $29$ phút, thời gian đi của ông thuộc nhóm $[30;33)$, do đó các giá trị thuộc nhóm $[30; 33)$ là các giá trị ngoại lệ, sau khi loại bỏ các giá trị ngoại lệ, ta có bảng thống kê sau đây
		\begin{center}
			\begin{tabular}{|c|c|c|c|c|c|}
				\hline
				Thời gian (phút) & {$[15; 18)$} & {$[18; 21)$} & {$[21; 24)$} & {$[24; 27)$} & {$[27; 30)$} \\
				\hline
				Số lượt & 22 & 38 & 27 & 8 & 4 \\
				\hline
			\end{tabular}
		\end{center}
		Khoảng biến thiên của mẫu số liệu khi chưa loại bỏ các giá trị ngoại lệ là $R=33-15=18$.\\
		Khoảng biến thiên của mẫu số liệu sau khi loại bỏ các giá trị ngoại lệ là $R=30-15=15$.\\
		Sau khi loại bỏ các giá trị ngoại lệ, mẫu số liệu mới có độ phân tán thấp hơn.\\
		Xét về khoảng tứ phân vị, ta có cỡ mẫu $n=99$.\\
		Gọi $x_1; x_2; \ldots; x_{99}$ là mẫu số liệu gốc gồm thời gian 99 lần đi xe buýt của ông Thắng.\\
		Ta có
		\begin{itemize}
			\item $x_1, \ldots, x_{22} \in[15; 18)$;
			\item $x_{23}, \ldots, x_{60} \in[18; 21)$; 
			\item $x_{61}, \ldots, x_{87} \in[21; 24)$; 
			\item $x_{88}, \ldots,x_{95} \in[24; 27)$; 
			\item $x_{96}, \ldots,x_{99} \in[27; 30)$.
		\end{itemize}
		Tứ phân vị thứ nhất của mẫu số liệu gốc là $x_{25} \in[18; 21)$. Do đó, tứ phân vị thứ nhất của mẫu số liệu ghép nhóm là \[Q_1=18+\dfrac{\dfrac{99}{4}-22}{38} \cdot(21-18)=\dfrac{2\,769}{152}.\]
		Tứ phân vị thứ ba của mẫu số liệu gốc là $x_{75} \in[21; 24)$. Do đó, tứ phân vị thứ ba của mẫu số liệu ghép nhóm là \[Q_3=21+\dfrac{\dfrac{3 \cdot 99}{4}-(22+38)}{27} \cdot(24-21)=\dfrac{271}{12}.\]
		Vậy khoảng tứ phân vị của mẫu số liệu ghép nhóm sau khi loại bỏ các giá trị ngoại lệ vẫn là $\Delta_Q=Q_3-Q_1=\dfrac{1\,991}{456}\approx 4{,}37$. 
	}
\end{ex}
%%%==============EX_3============%%%
\begin{ex}%[2D3V1-3]
	Cho bảng biểu diễn mẫu số liệu ghép nhóm về chiều cao của $42$ mẫu cây ở một vườn thực vật (đơn vị centimet). Tính khoảng tứ phân vị của mẫu số liệu ghép nhóm đó (làm tròn kết quả đến hàng phần mười).
	\begin{center}
		\begin{tabular}{|c|c|c|}
			\hline
			Nhóm & Tần số & Tần số tích luȳ \\
			\hline
			$[40; 45)$ & 5 & 5 \\
			{$[45; 50)$} & 10 & 15 \\
			{$[50; 55)$} & 7 & 22 \\
			{$[55; 60)$} & 9 & 31 \\
			{$[60; 65)$} & 7 & 38 \\
			{$[65; 70)$} & 4 & 42 \\
			\hline
			& $n=42$ & \\
			\hline
		\end{tabular}
	\end{center}
	\shortans[]{$12{,}7$}
	\loigiai{
		Số phần tử của mẫu là $n=42$.
		\begin{itemize}
			\item Ta có $\dfrac{n}{4}=\dfrac{42}{4}=10,5$ mà $5< 10,5< 15$.\\ Suy ra nhóm $2$ là nhóm đầu tiên có tần số tích lūy lớn hơn hoặc bằng $10{,}5$. 
			\item Xét nhóm $2$ là nhóm $[45; 50)$ có $s=45$; $h=5$; $n_2=10$ và nhóm 1 là nhóm $[40; 45)$ có $c f_1=5$. Áp dụng công thức, ta có tứ phân vị thứ nhất là \[Q_1=45+\left(\dfrac{10,5-5}{10}\right) \cdot 5=47{,}75.\]
			\item Ta có $\dfrac{3n}{4}=\dfrac{3\cdot 42}{4}=31,5$; mà $31< 31,5< 38$. \\
			Suy ra nhóm $5$ là nhóm đầu tiên có tần số tích lūy lớn hơn hoặc bằng $31{,}5$. 
			\item Xét nhóm $5$ là nhóm $[60; 65)$ có $t=60$, $l=5$; $n_5=7$ và nhóm 4 là nhóm $[55; 60)$ có $c f_4=31$.\\
			Áp dụng công thức, ta có tứ phân vị thứ ba là \[Q_3=60+\left(\dfrac{31{,}5-31}{7}\right) \cdot 5 \approx 60{,}4\,(\mathrm{cm}).\]
			Vậy khoảng tứ phân vị của mẫu số liệu ghép nhóm đā cho là
			\[\Delta_Q=Q_3-Q_1 \approx 60,4-47,75 \approx 12{,}7\,(\mathrm{cm})\]
		\end{itemize}
	}
\end{ex}
%%%=============EX_4=============%%%
\begin{ex}%[2D3H2-2]
	Để đánh giá chất lượng một loại pin điện thoại mới, người ta ghi lại thời gian nghe nhạc liên tục của điện thoại được sạc đầy pin cho đến khi hết pin cho kết quả sau
	\begin{center}
		\begin{tabular}{|c|c|c|c|c|c|} 
			\hline 
			Thời gian (giờ) & $\left[5;\,5{,}5\right)$ & $\left[5{,}5;\,6\right)$ & $\left[6;6{,}5\right)$ & $\left[6{,}5;\,7\right)$ & $\left[7;\,7{,}5\right)$ \\ 
			\hline 
			Số chiếc điện thoại (tần số) & $2$ & $8$ & $15$ & $10$& $5$\\ 
			\hline 
		\end{tabular}		
	\end{center}
	Tính phương sai của mẫu số liệu ghép nhóm trên (kết quả làm tròn đến hàng phần trăm).
	\shortans[]{$0{,}53$}
	\loigiai{
		Khoảng biến thiên $R=7{,}5-5=2{,}5$.\\
		Cỡ mẫu là $n=2+8+15+10+5=40$.\\
		Chọn giá trị đại diện cho mẫu số liệu ta có
		
		Thời gian trung bình là
		\[\overline{x}=\dfrac{5{,}25\cdot 2+5{,}75\cdot 8+15\cdot 6{,}25+10\cdot 6{,}75+5\cdot 7{,}25}{40}=6{,}35.
		\]
		Phương sai và độ lệch chuẩn là
		\[s^2=\dfrac{5{,}25^2\cdot 2+5{,}75^2\cdot 8+15\cdot 6{,}25^2+10\cdot 6{,}75^2+5\cdot 7{,}25^2}{40}-6{,}35^2=0{,}28.
		\]
		Suy ra $s=\sqrt{0{,}28} \approx 0{,}53$.
	}
\end{ex}

%%%=============EX_5=============%%%
\begin{ex}%[2D3V2-2]
	Bảng $1$, Bảng $2$ lần lượt biểu diễn mẫu số liệu ghép nhóm về nhiệt độ không khí trung bình các tháng năm $2021$ tại Hà Nội và Huế (đơn vị: độ)
	\begin{center}
		\begin{minipage}[t]{0.45\textwidth}
			\centering
			{\it Bảng 1}\\
			\vspace{0.2cm}
			\begin{tabular}{|c|c|c|}
				\hline
				Nhóm & Giá trị đại diện & Tần số \\
				\hline
				$[16{,}8; 19{,}8)$ & $18{,}3$ & 2 \\
				{$[19{,}8; 22{,}8)$} & $21{,}3$ & 3 \\
				{$[22{,}8; 25{,}8)$} & $24{,}3$ & 2 \\
				{$[25{,}8; 28{,}8)$} & $27{,}3$ & 1 \\
				{$[28{,}8; 31{,}8)$} & $30{,}3$ & 4 \\
				\hline
				& &$n=12$ \\
				\hline
			\end{tabular}
		\end{minipage}			
		\hspace{0.05\textwidth}
		\begin{minipage}[t]{0.45\textwidth}
			\centering
			{\it Bảng 2}\\
			\vspace{0.2cm}
			\begin{tabular}{|c|c|c|}
				\hline
				Nhóm & Giá trị đại diện & Tần số \\
				\hline
				$[16{,}8; 19{,}8)$ & $18{,}3$ & 1 \\
				{$[19{,}8; 22{,}8)$} & $21{,}3$ & 2 \\
				{$[22{,}8; 25{,}8)$} & $24{,}3$ & 3 \\
				{$[25{,}8; 28{,}8)$} & $27{,}3$ & 2 \\
				{$[28{,}8; 31{,}8)$} & $30{,}3$ & 4 \\
				\hline
				& &$n=12$ \\
				\hline			
			\end{tabular}
		\end{minipage}
	\end{center}
	
	\begin{flushright}
		\textit{(Nguồn: Niên giám thống kê $2021$, NXB Thống kê, $2022$)}
	\end{flushright}
	Gọi $s$, $s'$ lần lượt là độ lệch chuẩn của mẫu số liệu ghép nhóm của Hà Nội và Huế. Tính giá trị của biểu thức $T=s-s'$ (kết quả làm tròn đến hàng phần trăm).
	\shortans[]{$0{,}59$}
	\loigiai{
		\begin{itemize}
			\item Tại Hà Nội
			\begin{itemize}
				\item Từ Bảng $1$ ta có bảng thống kê sau
				\begin{center}
					\begin{tabular}{|c|c|c|}
						\hline
						Nhóm & Tần số & Tần số tích lũy\\
						\hline
						$[16{,}8; 19{,}8)$ & $2$ & 2 \\
						{$[19{,}8; 22{,}8)$} & $3$ & 5 \\
						{$[22{,}8; 25{,}8)$} & $2$ & 7 \\
						{$[25{,}8; 28{,}8)$} & $1$ & 8 \\
						{$[28{,}8; 31{,}8)$} & $4$ & 12 \\
						\hline
						& $n=12$& \\
						\hline
					\end{tabular}
				\end{center}
				\item Số phần tử của mẫu là $n=12$.
				\item Nhiệt độ trung bình của mẫu số liệu ghép nhóm được cho bởi Bảng $1$ là
				\[\overline{x}=\dfrac{2\cdot 18{,}3+3\cdot 21{,}3+2\cdot 24{,}3+1\cdot 27{,}3+4\cdot 30{,}3}{12}=\dfrac{297{,}6}{12}=24{,}8^\circ.
				\]
				Vậy phương sai biểu thị nhiệt độ của mẫu số liệu ghép nhóm được cho bởi Bảng $1$ là
				\[s^2=\dfrac{1}{12} \cdot \left[2\cdot (18{,}3-24{,}8)^2+3\cdot (21{,}3-24{,}8)^2+2\cdot (24{,}3-24{,}8)^2 \right.
				\]
				
				\[\left.+1\cdot (27{,}3-24{,}8)^2+4\cdot (30{,}3-24{,}8)^2 \right]=\dfrac{249}{12}=20{,}75.
				\]
				Độ lệch chuẩn biểu thị nhiệt độ của mẫu số liệu ghép nhóm trên là $s=\sqrt{20{,}75} \approx 4{,}56$.
			\end{itemize}
			\item Tại Huế
			\begin{itemize}
				\item Từ Bảng $2$ ta có bảng thống kê sau
				\begin{center}
					\begin{tabular}{|c|c|c|}
						\hline
						Nhóm & Giá trị đại diện & Tần số \\
						\hline
						$[16{,}8; 19{,}8)$ & $1$ & 1 \\
						{$[19{,}8; 22{,}8)$} & $2$ & 3 \\
						{$[22{,}8; 25{,}8)$} & $3$ & 6 \\
						{$[25{,}8; 28{,}8)$} & $2$ & 8 \\
						{$[28{,}8; 31{,}8)$} & $4$ & 12 \\
						\hline
						& $n=12$& \\
						\hline
					\end{tabular}
				\end{center}
				\item Số phần tử của mẫu là $n=12$.
				\item Nhiệt độ trung bình của mẫu số liệu ghép nhóm được cho bởi Bảng $2$ là
				\[\overline{x'}=\dfrac{1\cdot 18{,}3+2\cdot 21{,}3+3\cdot 24{,}3+2\cdot 27{,}3+4\cdot 30{,}3}{12}=\dfrac{309{,}6}{12}=25{,}8^\circ.
				\]
				Vậy phương sai biểu thị nhiệt độ của mẫu số liệu ghép nhóm được cho bởi Bảng $2$ là
				\[s'^2=\dfrac{1}{12} \cdot \left[1\cdot (18{,}3-25{,}8)^2+2\cdot (21{,}3-25{,}8)^2+3\cdot (24{,}3-25{,}8)^2 \right.
				\]
				
				\[\left.+2\cdot (27{,}3-25{,}8)^2+4\cdot (30{,}3-25{,}8)^2 \right]=\dfrac{189}{12}=15{,}75.
				\]
				Độ lệch chuẩn biểu thị nhiệt độ của mẫu số liệu ghép nhóm trên là $s'=\sqrt{15{,}75} \approx 3{,}97$.\\
				Vậy giá trị của biểu thức $T=s-s'=0{,}59$
			\end{itemize}
		\end{itemize}
		
	}
\end{ex}

%%%=============EX_6=============%%%
\begin{ex}%[2D3V1-3]
	Bảng  $3$ thống kê độ ẩm không khí trung bình các tháng năm $2021$ tại Đà Lạt và Vũng Tàu (đơn vị: $\%$). 
	\begin{center}
		\begin{tabular}{*{13}{|c}|}
			\hline
			Tháng &$1$&$2$&$3$&$4$&$5$&$6$&$7$&$8$&$9$&$10$&$11$&$12$ \\
			\hline
			Đà Lạt &83&79&79&87&87&87&88&89&90&91&88&86 \\
			\hline
			Vũng Tàu &75&77&78&77&79&79&81&79&81&83&80&77 \\
			\hline
		\end{tabular}\\
		\centering{\bf Bảng 3.}
	\end{center}
	\begin{flushright}
		\textit{(Nguồn: Niên giám thống kê $2021$, NXB Thống kê, $2022$)}
	\end{flushright}
	Gọi $\Delta_Q$, $\Delta'_Q$ lần lượt là khoảng tứ phân vị của mẫu số liệu ghép nhóm của Đà Lạt và Vũng Tàu. Tính giá trị biểu thức $K=\Delta '_Q-\Delta_Q$ (kết quả làm tròn đến hàng phần trăm).
	\shortans[]{$0{,}72$}
	\loigiai{
		\begin{itemize}
			\item Tại Đà Lạt
			\begin{itemize}
				\item Khoảng biến thiên của mẫu số liệu ghép nhóm của Đà Lạt là
				\[R=91{,}5-78{,}3=13{,}2 (\%).\]
				\item Từ bảng thống kê trên, ta có bảng thống kê của mẫu số liệu ghép nhóm của Đà Lạt
				
				\item Số phần tử của mẫu là $n=12$.
				\item Ta có $\dfrac{n}{4}=\dfrac{12}{4}=3$ mà $2< 3$. Suy ra nhóm $2$ là nhóm đầu tiên có tần số tích lũy lớn hơn hoặc bằng $3$. 
				\item Xét nhóm $2$ là nhóm $\left[81{,}6;\,84{,}9\right)$ có $s=81{,}6$; $h=3{,}3$; $n_2=1$ và nhóm $1$ là nhóm $\left[78{,}3;\,81{,}6\right)$ có $cf_1=2$.
				\item Áp dụng công thức, ta có tứ phân vị thứ nhất là
				\[Q_1=81{,}6+\left(\dfrac{3-2}{1} \right)\cdot 3{,}3=84{,}9\left(\% \right).
				\]
				\item Ta có $\dfrac{3n}{4}=\dfrac{3\cdot 12}{4}=9$ mà $3< 9< 10$. Suy ra nhóm $3$ là nhóm đầu tiên có tần số tích lũy lớn hơn hoặc bằng $9$.
				\item Xét nhóm $3$ là nhóm $\left[84{,}9;\,88{,}2\right)$ có $t=84{,}9$; $1=3{,}3$; $n_3=7$ và nhóm $2$ là nhóm $[81{,}6$; $84,9)$ có $cf_2=3$.
				\item Áp dụng công thức, ta có tứ phân vị thứ ba là
				\[Q_3=84{,}9+\left(\dfrac{9-3}{7} \right)\cdot 3{,}3\approx 87{,}7\left(\% \right).
				\]
				Vậy khoảng tứ phân vị của mẫu số liệu ghép nhóm của Đà Lạt là
				\[\Delta _{\mathbb{Q}}=Q_3-Q_1=87{,}7-84{,}9=2{,}8\left(\% \right).
				\]
			\end{itemize}			
			\item Tại Vũng Tàu
			\begin{itemize}
				\item Khoảng biến thiên của mẫu số liệu ghép nhóm của Vũng Tàu là \[R'=84{,}9-75=9{,}9 (\%).\]
				\item Từ bảng thống kê trên, ta có bảng thống kê của mẫu số liệu ghép nhóm của Vũng Tàu
				
				\item Số phần tử của mẫu là $n=12$.
				\item Ta có $\dfrac{n}{4}=\dfrac{12}{4}=3$ mà $5> 3$. Suy ra nhóm $1$ là nhóm đầu tiên có tần số tích lũy lớn hơn hoặc bằng $3$.
				\item Xét nhóm $1$ là nhóm $\left[75;\,78{,}3\right)$ có $s=75$; $h=3{,}3$; $n_1=5$.
				\item Áp dụng công thức, ta có tứ phân vị thứ nhất là
				\[Q_1^{'}=75+\dfrac{3}{5} \cdot 3{,}3=76{,}98\left(\% \right).
				\]
				\item Ta có $\dfrac{3n}{4}=\dfrac{3\cdot 12}{4}=9$ mà $5< 9< 11$. Suy ra nhóm $2$ là nhóm đầu tiên có tần số tích lũy lớn hơn hoặc bằng $9$.
				\item Xét nhóm $2$ là nhóm $\left[78{,}3;\,81{,}6\right)$ có $t=78{,}3$; $1=3{,}3$; $n_2=6$ và nhóm $1$ là nhóm $[75;\,78{,}3$) có $cf_1=5$.
				\item Áp dụng công thức, ta có tứ phân vị thứ ba là
				\[Q_3^{'}=78{,}3+\left(\dfrac{9-5}{6} \right)\cdot 3,3=80{,}5\left(\% \right).
				\]
				Vậy khoảng tứ phân vị của mẫu số liệu ghép nhóm của Vũng Tàu là
				\[\Delta _{\mathbb{Q}}^{'}=Q_3^{'}-Q_1^{'}=80{,}5-76{,}98=3{,}52\left(\% \right).
				\]
			\end{itemize}	
			Vậy giá trị biểu thức $K=\Delta '_Q-\Delta_Q=0{,}72\%$.
		\end{itemize}
	}
\end{ex}
\Closesolutionfile{ans}
\indapan{6}{ans/ans-De1-KQ}